\section{Disponibilidad de las observaciones}
\label{sec:disponibilidad}

La disponibilidad se evalúa para cada instrumento y mes del periodo analizado.
Expresa la proporción de observaciones que pueden utilizarse de acuerdo con los
criterios definidos para cada producto. No equivale al tiempo de funcionamiento
del instrumento ni permite interpretar los intervalos sin datos como ausencia
de precipitación.

El calendario de la Figura~\ref{fig:calendario-disponibilidad} ofrece una
visión conjunta de la cobertura temporal. La Tabla~\ref{tab:disponibilidad-mensual}
recoge los valores mensuales y permite distinguir entre meses no analizados,
meses analizados sin observaciones válidas y meses con cobertura parcial o
completa.

\IfFileExists{figuras/calendario_disponibilidad.png}{%
  \begin{figure}[htbp]
    \centering
    \includegraphics[width=\textwidth]{calendario_disponibilidad.png}
    \caption{Disponibilidad temporal de las observaciones por instrumento
      durante el periodo analizado. Los intervalos sin registros no se
      interpretan como ausencia de precipitación.}
    \label{fig:calendario-disponibilidad}
  \end{figure}
}{%
  \begin{figure}[htbp]
    \centering
    \fbox{\parbox{0.88\textwidth}{\centering
      Calendario conjunto pendiente de generación a partir de los resultados
      preparados del periodo.}}
    \caption{Disponibilidad temporal por instrumento. Producto pendiente.}
    \label{fig:calendario-disponibilidad}
  \end{figure}
}

\subsection{Resumen mensual}
\label{sec:resumen-disponibilidad}

\IfFileExists{tablas/disponibilidad_mensual.tex}{%
  \begin{table}[htbp]
    \centering
    \caption{Resumen mensual de disponibilidad por instrumento.}
    \label{tab:disponibilidad-mensual}
    \input{tablas/disponibilidad_mensual}
  \end{table}
}{%
  \begin{table}[htbp]
    \centering
    \caption{Resumen mensual de disponibilidad por instrumento. Producto
      pendiente.}
    \label{tab:disponibilidad-mensual}
    \begin{tabular}{@{}ll@{}}
      \toprule
      Producto & Estado \\
      \midrule
      Tabla mensual común & Pendiente de generación \\
      \bottomrule
    \end{tabular}
  \end{table}
}

En la tabla, un guion indica que no existe un resultado preparado para ese mes
y producto. Un valor de 0\,\% indica que el mes sí está analizado, pero no contiene
observaciones válidas según el criterio correspondiente.

Ambos productos se generan a partir de los resultados mensuales preparados. Si
alguno no existe, el documento muestra un sustituto que identifica expresamente
el producto pendiente. No se extraen conclusiones sobre meses concretos sin
resultados preparados que las respalden.
